% !TeX program = lualatex
% Definition and sources only; this file is not an LLM solution attempt.
\documentclass[11pt,a4paper]{article}
\usepackage[margin=25mm]{geometry}
\usepackage{fontspec}
\setmainfont{lmroman10-regular.otf}[BoldFont=lmroman10-bold.otf,
 ItalicFont=lmroman10-italic.otf,BoldItalicFont=lmroman10-bolditalic.otf]
\usepackage{amsmath,amssymb,mathtools}
\usepackage{unicode-math}
\setmathfont{latinmodern-math.otf}
\AtBeginDocument{\let\setminus\smallsetminus}
\usepackage{xurl,hyperref}
\hypersetup{colorlinks=true,urlcolor=blue}
\setcounter{secnumdepth}{0}
\setlength{\parindent}{0pt}
\setlength{\parskip}{5pt}
\setlength{\emergencystretch}{3em}
\begin{document}
\section*{A class of Hamiltonian systems}\label{statement}

\subsection{Definitions and mathematical statement}
Let $n\ge1$ and $m>2n$ be integers, and let $H_{2n},H_m\in\mathbb R[x,y]$
be homogeneous polynomials of the indicated degrees. Set
\[
 H=H_{2n}+H_m,\qquad \dot x=H_y,\qquad \dot y=-H_x,
\]
and assume that the origin is a center: a punctured neighborhood is filled
by periodic orbits surrounding this isolated equilibrium. Its period
annulus is the maximal connected family of such surrounding periodic orbits.
On each regular level oval $\gamma_h$ in that annulus, let
$T(h)>0$ be the least positive period of the Hamiltonian flow.

A critical period means an isolated stationary point of this function:
$T'(h)=0$ at an interior energy value. Equivalently one may use a smooth
transverse parameter, since stationary points are unchanged by a regular
reparametrization. The question is whether there is at most one critical
period throughout the full period annulus of the origin.

The count concerns the actual Hamiltonian time, and a nonconstant time
rescaling is not allowed. Degenerate centers are included, so their periods
may diverge as the energy approaches the center. Endpoints corresponding to
the center or the outer boundary are not counted as interior critical
periods. The source proves the bound in the range $m\ge4n-2$; the remaining
range is $2n<m<4n-2$, beginning with a quartic-plus-quintic Hamiltonian.
These degree restrictions do not impose any extra sign assumption beyond
the stated existence of a center.
\subsection{Short English statement}
Does the period annulus of a center for a two-term homogeneous Hamiltonian have at most one critical period?
\subsection{Sources}
\begin{itemize}
\item[S1] ULAM.AI and the UnsolvedMath Contributors, supplied problems.json, catalogue entry, AMR Open Problem Lists. Catalogue identity and source transcription; CC BY 4.0.
\url{https://huggingface.co/datasets/ulamai/UnsolvedMath}.
\item[S2] Gasull - Some open problems in low dimensional dynamical systems (2020), Problem environment 12: A class of Hamiltonian systems. 
\url{https://arxiv.org/abs/2012.02524}.
\end{itemize}
\end{document}
